\chapter{Survey Materials} \label{sec:SurveyMaterials}

Below is a collection of various materials used for the surveys conducted to evaluate
Praxsmith.

\section{Outreach Post}

Hi Everyone!

Want to participate in cutting edge game development research? Read on for more
information!

My name is Jacob Kelleran, and I'm a Computer Science Master's student at Cal Poly, San Luis
Obispo in California. I'm writing to you today to ask for your assistance with some research I'm
conducting on Emergent Narrative Game Development. You'll effectively be writing character
dialogue interactions with a new framework we are developing.

This study requires filling out a short survey, then using a new tool to expand a short narrative
experience, then filling out another survey, and finally having a discussion with the authors of the
tool to evaluate its effectiveness. All parts of this process can be done remotely, and contact
will be done through email. This study is investigating the validity, effectiveness, and efficiency of
the new tool being developed. This project has been approved by the Cal Poly San Luis Obispo
IRB (IRB Protcol \#2026-080).

To be eligible for the study, you must be 18 years of age or older and be able to either come to Cal
Poly for the discussion or use Zoom. The study takes approximately at least 3 but no more than 6
hours over the course of a week regardless of mode.

Please message me if you would like to participate. I will provide you with a link to our pre-
experiment survey that includes the full consent information and study details.

The gaming industry is the largest entertainment industry and grows more accessible every year.
Emergent narratives are a notoriously challenging task, and by participating in this study, you will
be a great help to furthering this research!

Thank you for your time!

\section{Initial Email}

Hello!

I am reaching out to you because you completed the pre-experiment survey for my
masters thesis. Here is a link to the project that you will be experimenting with:

https://praxsmith.emuman.net/

Using the custom-built Praxsmith language, you can create various worlds within
the provided sandbox. You will define actors, their goals, and the relations and
social practices that bind them together. You can then have them autonomously
make decisions through the planning system or manually intervene in their place.
The project is still in development, so expect some bugs and rough edges here and
there. There is thorough documentation on the language itself on the website, as
well as a few examples you can build off of. You will probably want to keep the
documentation open on the side for quick reference during your development.

Feel free to play around with it in whatever ways you like, even if you're
just messing around with an existing demo or creating a small one of your own.
You may spend as little or as much time on this as you'd like, but try to get
familiar with the systems and make something that you find interesting.

After you are done experimenting, fill out the post-experiment form to let us
know how your experience went.

If you have any questions or run into any issues, please don't hesitate to reach out.
We will be available to fix bugs, clarify any confusion, and offer advice as quickly
as possible throughout the experience.

Thank you for your participation!

\section{Pre-Experience Survey Questions}

Likert scale (1-5):

\begin{itemize}
    \item How much do you play games?
    \item Rate your experience developing games.
    \item What level of confidence do you have writing stories and characters?
    \item How familiar are you with programming or scripting languages?
    \item How often do you engage with narrative media outside of games? (books, films, interactive fiction, etc.)
    \item How much experience do you have with worldbuilding or collaborative storytelling? (e.g. TTRPGs, fanfiction, improv)
\end{itemize}

What genres of games do you enjoy playing? (select all that apply)

\begin{itemize}
    \item Action / Action-Adventure
    \item Shooter (FPS / TPS)
    \item Fighting
    \item Platformer
    \item Puzzle
    \item Strategy / Tactics (e.g. Civilization, XCOM, StarCraft)
    \item Simulation / Management (e.g. RimWorld, Cities: Skylines)
    \item Sports / Racing
    \item Horror / Survival
    \item Sandbox / Open World
    \item Life Simulator (e.g. The Sims, Animal Crossing)
    \item Role-Playing Game (RPG)
    \item Japanese RPG (JRPG)
    \item Tabletop RPG (e.g. D\&D, Pathfinder)
    \item Visual Novel
    \item Interactive Fiction / Text Adventure
    \item Walking Simulator / Narrative Exploration
    \item Social Deduction (e.g. Among Us, Town of Salem)
    \item Mobile / Casual
    \item I don't play video games
\end{itemize}

\section{Post-Experience Survey Questions}

Likert scale (1-5):

\begin{itemize}
    \item How easy were the rules of the language to understand?
    \item How easy was the language to use?
    \item How natural did the language feel to write in?
    \item How well did the language support the story you had in mind?
    \item How satisfied were you with the resulting game?
    \item Was there sufficient documentation for you to find your way around?
    \item Was the template effective in helping you get started?
\end{itemize}

What did you find most effective or enjoyable about working with the language?

Were there any issues you encountered writing in the new language?

Was there anything you wish you were able to do but couldn’t figure out how?

Was there anything you wish you were able to put in your final product but didn’t have time?

Was there anything that you found yourself repeating often and wished you didn’t have to?

Was there anything about the language that surprised you, positively or negatively?

Is there anything you would change about the language or template if you could?
